% ==============================================================================
% Fichier : front/02_mode_emploi.tex
% Titre   : Mode d'emploi — TPIN
% ==============================================================================

\chapter*{Mode d'emploi de ce cours}
\addcontentsline{toc}{chapter}{Mode d'emploi}
\markboth{Mode d'emploi}{Mode d'emploi}

\section*{Structure générale}

Ce livre correspond à l'UE \textbf{SEC4052}. Il couvre les \textbf{parties IV à VI}
(chapitres 14 à 32)~: forensique opérationnelle, anti-forensique et défis
émergents, crypto-forensique. Il suppose acquis le prologue et les parties I
à III (chapitres 1 à 13), enseignés dans l'UE SEC4031, \emph{Introduction à
l'Investigation Numérique}~; le tableau ci-dessous rappelle l'ensemble du
parcours.

\begin{tableaucours}{lp{3.8cm}p{3.8cm}c}{%
    \tableauHeader{Partie} &
    \tableauHeader{Thème} &
    \tableauHeader{Compétences visées} &
    \tableauHeader{Niveau}
}
0 & Prologue opérationnel & Appréhender la réalité du terrain & \niveaubase \\
I & Fondements épistémologiques et historiques & Comprendre l'origine et les bases théoriques & \niveaubase \\
II & Processus et méthodologies & Conduire une investigation structurée & \niveaubase\,--\,\niveaumoyen \\
III & Cadre normatif et judiciaire & Opérer dans la légalité & \niveaubase \\
IV & Forensique opérationnelle & Analyser système, réseau, mobile, cloud & \niveaumoyen\,--\,\niveauexpert \\
V & Anti-forensique et défis émergents & Contrer les techniques d'occultation & \niveauexpert \\
VI & Crypto-forensique et contributions & Maîtriser les fondements post-quantiques & \niveauexpert \\
\end{tableaucours}

\section*{Les environnements visuels}

\begin{boxdef}
\textbf{Définition}. Introduit un terme technique avec sa définition formelle.
Tous les termes définis sont indexés en fin d'ouvrage.
\end{boxdef}

\begin{boxobjectifs}
Objectifs du chapitre. Liste les compétences mesurables attendues à l'issue
du chapitre. Ces objectifs correspondent aux critères d'évaluation.
\end{boxobjectifs}

\begin{boxCRO}
Analyse CRO. Évalue systématiquement une technique ou un protocole selon
les trois dimensions du Trilemme : Confidentialité~($\mathcal{C}$),
Fiabilité~($\mathcal{R}$), Opposabilité~($\mathcal{O}$).
\end{boxCRO}

\begin{boxscenario}{Titre du scénario}
Scénario opérationnel. Présente un cas réel ou fictif ancré dans la réalité
camerounaise ou africaine. Le cas MBONGO (fraude financière 12M~FCFA) sert
de fil rouge depuis l'UE SEC4031 et reparaît dans la Partie~IV.
\end{boxscenario}

\begin{boxPNC}
Application OPJ/DGSN. Traduit le concept académique en procédure opérationnelle
destinée aux Officiers de Police Judiciaire.
\end{boxPNC}

\begin{boxcustody}
Chaîne de Custody. Détaille les étapes et formulaires nécessaires à la
préservation de la valeur probante d'une preuve numérique.
\end{boxcustody}

\begin{boxartefact}{Nom de l'artefact}
Artefact forensique. Décrit un artefact numérique spécifique : sa localisation
sur le système, la valeur qu'il porte pour l'investigation, et sa fragilité.
\end{boxartefact}

\begin{boxjuris}
Jurisprudence \& Droit applicable. Référence les textes de loi, articles et
jurisprudences pertinents, avec accent sur le droit camerounais et africain.
\end{boxjuris}

\begin{boxalert}
Attention. signale un risque procédural, légal ou technique pouvant
invalider une investigation.
\end{boxalert}

\begin{boxresume}
À retenir. Synthèse des points essentiels du chapitre.
\end{boxresume}

\begin{boxexo}
Exercice. Travaux pratiques gradués. Chaque exercice porte un indicateur
de niveau : \niveaubase\ fondamental, \niveaumoyen\ intermédiaire,
\niveauexpert\ expert.
\end{boxexo}

\section*{Conventions typographiques}

\begin{itemize}
    \item Les \termecle{termes clés} apparaissent en \textbf{gras bleu marine}
          et sont indexés automatiquement.
    \item Les \sigle{ACRONYMES} utilisent des petites capitales.
    \item Les références aux normes utilisent la commande \norme{ISO/IEC~27037}.
    \item Les références aux lois camerounaises utilisent \loicam{2010/012}.
    \item Les hachages cryptographiques sont formatés ainsi :
          \hash{SHA-256}{\texttt{3a4b5c6d...}}.
    \item L'impact CRO inline est noté : \cro{+0.8}{-0.3}{+0.6}.
    \item L'ordre de volatilité est indiqué : \volatilite{1}{RAM}.
\end{itemize}

\section*{Matériel d'accompagnement}

Ce cours est complété par :
\begin{itemize}
    \item \textbf{Images de pratique forensique}, images E01 préparées pour
          les TP Autopsy.
    \item \textbf{Captures réseau (.pcap)}, traces Wireshark pour les TP
          de forensique réseau.
    \item \textbf{Machines virtuelles}, VM Kali Linux et VM Windows avec
          artefacts préinstallés (clé USB fournie en TP).
    \item \textbf{Formulaires officiels}, formulaire de chaîne de custody
          (Annexe~C), adapté du formulaire type DGSN.
\end{itemize}
